% TC-0001 test-case definition (WSP test-case template).
\documentclass[letterpaper]{article}

\usepackage[letterpaper,margin=1in]{geometry}

% Compile from the repository root so this path resolves to the pinned WSP
% submodule.
\input{wsp/testing/test-case-library.tex}

\begin{document}

\def\TCID{TC-0001}
\def\TCTitle{Member table of the five classes}
\def\TCPurpose{%
  Verify that the member table lists the five classes of MSComCtl2 in order, that every class lists named members of a known kind, and that every event is listed as a method.}
\def\TCPriority{\TCPriorityHigh}

\def\TCRequirementRef{REQ-0006}
\def\TCDesignRef{Not applicable.}
\def\TCFeatureRef{\texttt{src/mscomctl2\_tables.cpp}; \texttt{include/wfc/mscomctl2/mscomctl2.hpp}; \texttt{tests/mscomctl2\_tables\_test.cpp}.}
\def\TCTestTechnique{%
  Structural check of generated data: every class, member and event row is visited.}

\def\TCPreconditions{%
  The library is configured and built with \texttt{cmake -S . -B out/build} and \texttt{cmake --build out/build}.}
\def\TCEnvironment{%
  Windows (CI: x86, x64 and ARM64, static and shared build; Linux Winelib and WebAssembly are experimental). CMake 3.28 or later, CTest, a C++20 compiler; Debug configuration locally (x64).}
\def\TCAssumptions{%
  Not applicable.}

\def\TCInputData{%
  The table returned by \texttt{wfc::ui::mscomctl2\_classes()}.}
\def\TCInitialState{%
  No output of an earlier run is used. CTest test name: \texttt{TC-0001-tables}.}

\def\TCProcedure{%
  \begin{enumerate}
    \item Run \texttt{mscomctl2\_tables\_test} through CTest.
    \item The program compares the number and the names of the classes with \texttt{Animation, UpDown, MonthView, DTPicker, FlatScrollBar} in that order.
    \item For every class it checks that the class has members, that every member has a non-empty name and a kind in \texttt{g, l, s, m}, and that every event has the kind \texttt{m}.
  \end{enumerate}}
\def\TCExpectedResult{%
  The program exits with status 0 and prints \texttt{mscomctl2: 5 classes}. Any violation prints a line starting with \texttt{FAIL:} and gives a nonzero exit status.}
\def\TCPostConditions{%
  No files are written.}

\TCTemplate

\end{document}
